% !TeX program = pdflatex
%==============================================================================
%  Avaliação Teórica 02 --- Aulas 07 a 10 (Bloco II, classificação)
%
%  O gabarito sai de "Avaliacao teorica 02 - gabarito.tex", que apenas liga a
%  opção [gabarito] e inclui este arquivo --- fonte única, dois PDFs.
%
%  AS TRÊS QUESTÕES FORAM SORTEADAS da "Lista de revisao 07-10", com cobertura
%  (uma aula por questão, três das quatro aulas). O sorteio é reproduzível:
%
%      import numpy as np
%      rng = np.random.default_rng(20261112)      # a data, aaaammdd
%      aulas = sorted(rng.choice([7, 8, 9, 10], size=3, replace=False).tolist())
%      exercicios = {7: [1, 2], 8: [3, 4], 9: [5, 6], 10: [7, 8, 9]}
%      sorteados = [int(rng.choice(exercicios[a])) for a in aulas]
%      # -> aulas [8, 9, 10]; exercícios [3, 6, 9]
%
%  Resultado:
%      Aula 08 -> Exercício 3  ex-03-sorteio-classificadores (sortear entre
%                              dois classificadores)
%      Aula 09 -> Exercício 6  ex-06-kernel-radial (o kernel radial em
%                              qualquer dimensão)
%      Aula 10 -> Exercício 9  ex-09-gini-entropia (quando o Gini e a entropia
%                              discordam)
%
%  O sorteio rodou uma vez e vale. O algoritmo é o da Avaliação Teórica 01: com
%  a semente 20260914 e a lista 01-06 de então, ele dá as aulas [2, 3, 6] e os
%  Ex. 5, 7 e 13 registrados lá.
%
%  NÃO HÁ TEXTO DE EXERCÍCIO NESTE ARQUIVO. Os enunciados e as soluções vêm por
%  \input dos MESMOS arquivos que a lista de revisão usa, em exercicios-07-10/.
%  Aqui ficam só o cabeçalho da prova, com a notação da caixa de instruções, as
%  questões e os critérios de correção, que são específicos da avaliação.
%
%  O \provatrue abaixo liga a pontuação por item: o comando \pts, escrito dentro
%  dos arquivos de exercício, some na lista e aparece aqui.
%
%  A PROVA NÃO REVELA DE QUE AULA VEM CADA QUESTÃO: o \questao não leva marcador
%  de aula, e os títulos ficam vazios porque os exercícios da lista 07-10 não
%  têm título. A notação de classificação entra como argumento opcional do
%  \cabecalhoprova; sem ele, a caixa imprimiria a de regressão, da prova 01.
%==============================================================================
\providecommand{\gabaritoopt}{}
\documentclass[11pt,a4paper]{article}
\usepackage[\gabaritoopt]{../recursos/latex/estilo-avaliacao}

\provatrue

\begin{document}
\cabecalhoprova[{$(\X,Y)$ é um vetor aleatório com $\X\in\R^p$ e $Y$ num conjunto finito
            $\mathcal{C}$ de classes; no caso binário, $\mathcal{C}=\{0,1\}$, a classe
            $Y=1$ é a positiva e $\eta(\x)=\Prob(Y=1\mid\X=\x)$;
            $\Dados=\{(\X_i,Y_i)\}_{i=1}^n$ é a amostra de treino, i.i.d.; a perda é a
            0--1, com risco $\risco(g)=\Prob\big(Y\ne g(\X)\big)$, salvo indicação em
            contrário.}]{02}{07 a 10}{Classificadores Gaussianos, Métricas, SVM, KNN e
Árvores de Classificação}{2 horas}{12,0 pontos}

%==============================================================================
\questao{}{4,0 pontos}
%==============================================================================
% Aula 08 --- classificadores sorteados: combinações convexas de pontos ROC.
Diego é analista antifraude num banco e tem dois sistemas que já foram avaliados, $A$ e
$B$. Tome a fraude como a classe positiva, $Y=1$. As taxas de falsos positivos e de
verdadeiros positivos são $(\mathrm{FPR},\mathrm{TPR})=(0{,}02;\ 0{,}50)$ para $A$ e
$(0{,}08;\ 0{,}80)$ para $B$. O banco aceita mandar para revisão manual no máximo $5\%$ das
transações legítimas, e nenhum dos dois sistemas aproveita esse limite: $A$ fica abaixo
dele e $B$ o ultrapassa. Diego então propõe sortear: para cada transação, uma moeda com
probabilidade $\lambda$ de dar cara, lançada independentemente de tudo o mais, escolhe se a
decisão será a de $A$ (cara) ou a de $B$ (coroa). Seja $g_\lambda$ esse classificador.
\begin{itens}
  \item \pts{1,6} Mostre que
        $\mathrm{FPR}(g_\lambda)=\lambda\,\mathrm{FPR}(A)+(1-\lambda)\,\mathrm{FPR}(B)$ e
        que o mesmo vale para a TPR, isto é, que o ponto de $g_\lambda$ no plano
        (FPR, TPR) está no segmento que liga os pontos de $A$ e de $B$. Determine $\lambda$
        para que a taxa de falsos positivos seja exatamente $0{,}05$ e calcule a taxa de
        verdadeiros positivos obtida. Que classificadores sorteados dão os pontos da
        diagonal?
  \item \pts{1,2} Um terceiro sistema, $C$, tem $(\mathrm{FPR},\mathrm{TPR})=(0{,}70;\ 0{,}40)$.
        Mostre que o classificador que inverte todas as decisões de $C$ tem o ponto
        $(0{,}30;\ 0{,}60)$ e, em geral, que inverter as decisões leva
        $(\mathrm{FPR},\mathrm{TPR})$ a $(1-\mathrm{FPR},1-\mathrm{TPR})$. Se um escore tem
        curva ROC abaixo da diagonal, com área $\mathrm{AUC}<1/2$, qual é a curva e qual é a
        área do escore com o sinal trocado?
  \item \pts{1,2} Com custos $\ell_0$ por falso negativo e $\ell_1$ por falso positivo, o
        custo esperado de um classificador $g$ é
        $\risco_\ell(g)=\ell_1\Prob(g(\X)=1,\,Y=0)+\ell_0\Prob(g(\X)=0,\,Y=1)$. Mostre que
        $\risco_\ell(g_\lambda)=\lambda\risco_\ell(A)+(1-\lambda)\risco_\ell(B)$ e conclua
        que, sem a restrição do banco, sortear nunca custa menos que o melhor dos dois
        sistemas. Para que serve, então, o sorteio?
\end{itens}

\begin{solucao}
\textbf{(a)} Seja $Z$ a moeda, com $\Prob(Z=1)=\lambda$, independente de $(\X,Y)$. Pela lei
da probabilidade total, condicionando em $Z$,
\begin{align*}
  \Prob\big(g_\lambda(\X)=1\mid Y=0\big)
  &=\lambda\,\Prob\big(g_A(\X)=1\mid Y=0,\,Z=1\big)\\
  &\quad+(1-\lambda)\,\Prob\big(g_B(\X)=1\mid Y=0,\,Z=0\big),
\end{align*}
em que se usou $\Prob(Z=1\mid Y=0)=\Prob(Z=1)=\lambda$, pela independência. Também pela
independência, condicionar em $Z$ não altera a distribuição de $\X$ dado $Y$, e as duas
probabilidades condicionais são as taxas de falsos positivos de $A$ e de $B$. Assim,
$\mathrm{FPR}(g_\lambda)=\lambda\,\mathrm{FPR}(A)+(1-\lambda)\,\mathrm{FPR}(B)$, e a mesma
conta, condicionando em $Y=1$, dá a TPR. O ponto de $g_\lambda$ é a combinação convexa dos
pontos de $A$ e de $B$ com pesos $\lambda$ e $1-\lambda$, e percorre o segmento entre eles
quando $\lambda$ vai de $0$ a $1$.

Para a taxa de falsos positivos valer $0{,}05$,
\[
  0{,}02\,\lambda+0{,}08\,(1-\lambda)=0{,}05\iff 0{,}06\,\lambda=0{,}03\iff\lambda=\tfrac12 ,
\]
e a taxa de verdadeiros positivos é $\tfrac12\cdot0{,}50+\tfrac12\cdot0{,}80=0{,}65$: Diego usa
o limite do banco inteiro e detecta $65\%$ das fraudes, contra $50\%$ de $A$. A diagonal
vem do sorteio entre os dois classificadores triviais, $g\equiv0$, no ponto $(0,0)$, e
$g\equiv1$, no ponto $(1,1)$: sortear $g=1$ com probabilidade $q$ dá o ponto $(q,q)$.

\smallskip
\textbf{(b)} Se $g'=1-g$, então $g'(\X)=1$ exatamente quando $g(\X)=0$, e
\[
  \mathrm{TPR}(g')=\Prob\big(g(\X)=0\mid Y=1\big)=1-\mathrm{TPR}(g),
  \qquad
  \mathrm{FPR}(g')=1-\mathrm{FPR}(g).
\]
Para $C$, o ponto invertido é $(1-0{,}70;\ 1-0{,}40)=(0{,}30;\ 0{,}60)$, acima da diagonal: um
sistema que erra sistematicamente carrega informação, desde que se faça o contrário do que
ele diz. A aplicação $(a,b)\mapsto(1-a,1-b)$ é a simetria em relação ao ponto
$(\tfrac12,\tfrac12)$. Trocar o sinal do escore troca ``escore acima do corte'' por
``escore abaixo do corte'', o que inverte as decisões em cada corte; a curva do escore com o
sinal trocado é, portanto, a simétrica da original em relação a $(\tfrac12,\tfrac12)$. Essa
simetria leva a região do quadrado unitário abaixo de uma curva na região acima da curva
simétrica e preserva áreas. A área sob a curva nova é, então, a área acima da curva antiga,
$1-\mathrm{AUC}>1/2$.

\smallskip
\textbf{(c)} Condicionando de novo na moeda, e usando a independência,
\[
  \Prob\big(g_\lambda(\X)=1,\,Y=0\big)
  =\lambda\,\Prob\big(g_A(\X)=1,\,Y=0\big)+(1-\lambda)\,\Prob\big(g_B(\X)=1,\,Y=0\big),
\]
e o mesmo para $\Prob(g_\lambda(\X)=0,\,Y=1)$. Multiplicando a primeira igualdade por
$\ell_1$, a segunda por $\ell_0$ e somando,
$\risco_\ell(g_\lambda)=\lambda\risco_\ell(A)+(1-\lambda)\risco_\ell(B)$. Uma média ponderada
de dois números nunca é menor que o menor deles, e o sorteio nunca custa menos que o melhor
dos dois sistemas: sem restrição, basta escolher um deles. O sorteio serve para atingir um
ponto que nenhum dos dois atinge, como a taxa de falsos positivos de exatamente $0{,}05$
permitida pelo banco: respeitando esse limite, $g_{1/2}$ detecta $65\%$ das fraudes, e $A$,
só $50\%$.
\end{solucao}


\begin{criterio}
\textbf{(a)} 0,8 pela fórmula da FPR: 0,4 por condicionar na moeda, com
$\Prob(Z=1\mid Y=0)=\lambda$, e 0,4 por justificar, pela independência, que as duas
probabilidades condicionais são as taxas de falsos positivos de $A$ e de $B$. Afirmar a
combinação convexa sem usar a independência vale 0,3 desses 0,8, e a TPR, que sai da
mesma conta, não precisa ser refeita. 0,5 pelos números: 0,3 por $\lambda=1/2$ e 0,2 pela
TPR de $0{,}65$. 0,3 pela diagonal: sortear entre $g\equiv0$ e $g\equiv1$, isto é, declarar
positivo com probabilidade $q$ sem olhar $\x$, dá o ponto $(q,q)$.

\textbf{(b)} 0,3 pelo ponto $(0{,}30;\ 0{,}60)$; 0,4 pelo caso geral,
$\mathrm{TPR}(1-g)=1-\mathrm{TPR}(g)$ e o mesmo para a FPR; 0,5 pelo escore com o sinal
trocado, 0,25 pela curva, a simétrica da original em relação a $(\tfrac12,\tfrac12)$, e
0,25 pela área, $1-\mathrm{AUC}>1/2$. Aceitar a área justificada pela simetria ou pela
área acima da curva original; não é exigido tratar empates entre escores.

\textbf{(c)} 0,5 pela linearidade do custo esperado, que repete a conta do item (a) com
as probabilidades conjuntas; remeter ao item (a) vale integralmente se ficar claro que
a conta se faz para os dois tipos de erro e que os custos entram como constantes. 0,3
por concluir que uma média ponderada nunca é menor que o menor dos dois custos. 0,4 pela
utilidade do sorteio: atingir uma taxa de falsos positivos que nenhum dos dois sistemas
atinge, como os $0{,}05$ permitidos pelo banco, com $65\%$ das fraudes detectadas contra
$50\%$ de $A$. Dizer apenas que o sorteio combina os dois sistemas, sem a restrição do
banco, vale 0,1 desses 0,4.
\end{criterio}

%==============================================================================
\questao{}{4,0 pontos}
%==============================================================================
% Aula 09 --- produtos e exponenciais de kernels; o kernel radial em qualquer dimensão.
Júlia estuda matemática aplicada e lê nas notas que o \emph{kernel} radial
$k(\x,\x')=\exp\!\big(-\gamma\norm{\x-\x'}^2\big)$, com $\gamma>0$, é um \emph{kernel} em
$\R^p$; o exemplo das notas, porém, só exibe a função $\varphi$ para $p=1$. Ela decide provar
o caso geral. Lembre que $k:\R^p\times\R^p\to\R$ é um \emph{kernel} se existem um espaço de
Hilbert $\mathcal H$ e uma função $\varphi:\R^p\to\mathcal H$ com
$k(\x,\x')=\inner{\varphi(\x)}{\varphi(\x')}_{\mathcal H}$ para todos $\x,\x'$, e que, pelo
teorema visto em aula, isso acontece se e somente se $k$ é simétrica e, para todo $n\ge1$ e
quaisquer $\x_1,\dots,\x_n\in\R^p$, a matriz $\big[k(\x_i,\x_j)\big]_{i,j=1}^n$ é positiva
semidefinida. Use à vontade que o produto interno $\inner{\x}{\x'}$ e as funções constantes
iguais a $c\ge0$ são \emph{kernels}, e que a soma de dois \emph{kernels} e o produto de um
\emph{kernel} por uma constante $c\ge0$ também são.
\begin{itens}
  \item \pts{1,2} Mostre que, se $k_1$ e $k_2$ são \emph{kernels}, o produto
        $k(\x,\x')=k_1(\x,\x')\,k_2(\x,\x')$ também é. (Escreva uma das duas matrizes pela
        decomposição espectral.)
  \item \pts{1,0} Mostre que, se $k$ é um \emph{kernel} e $h:\R^p\to\R$ é uma função
        qualquer, então $h(\x)\,h(\x')\,k(\x,\x')$ é um \emph{kernel}. Mostre também que
        $e^{k(\x,\x')}$ é um \emph{kernel}, usando a série da exponencial e o fato de que um
        limite de matrizes positivas semidefinidas é positivo semidefinido.
  \item \pts{0,8} Conclua que o \emph{kernel} radial é um \emph{kernel} em $\R^p$, para todo
        $p$ e todo $\gamma>0$. (Escreva $\norm{\x-\x'}^2$ em termos de $\norm{\x}^2$,
        $\norm{\x'}^2$ e $\inner{\x}{\x'}$.)
  \item \pts{1,0} Mostre que, para todo \emph{kernel},
        $\norm{\varphi(\x)-\varphi(\x')}_{\mathcal H}^2=k(\x,\x)-2\,k(\x,\x')+k(\x',\x')$, e
        calcule essa distância para o \emph{kernel} radial. Conclua que $\varphi$ leva $\R^p$
        na esfera unitária de $\mathcal H$ e que pontos distintos ficam, em $\mathcal H$, a
        distância menor que $\sqrt2$. Com $\gamma=0{,}5$ e com $\gamma=50$, a que distância
        dois pontos precisam estar em $\R^p$ para ficarem à distância $1$ em $\mathcal H$?
        Interprete o resultado em termos da flexibilidade da SVM. \emph{Dado:}
        $\sqrt{2\log2}\approx1{,}177$, com $\log$ o logaritmo natural.
\end{itens}

\begin{solucao}
\textbf{(a)} Sejam $\x_1,\dots,\x_n\in\R^p$, $A=\big[k_1(\x_i,\x_j)\big]$ e
$B=\big[k_2(\x_i,\x_j)\big]$, simétricas e positivas semidefinidas pelo teorema. A matriz de
$k$ nesses pontos é o produto entrada a entrada, $(A\circ B)_{ij}=A_{ij}B_{ij}$, que é
simétrica. Pelo teorema espectral, $B=\sum_{m=1}^n\lambda_m\bm v_m\bm v_m^\top$, com
autovalores $\lambda_m\ge0$, porque $B$ é positiva semidefinida, e autovetores ortonormais
$\bm v_m\in\R^n$; em coordenadas, $B_{ij}=\sum_m\lambda_m v_{mi}v_{mj}$. Para $\bm a\in\R^n$,
\begin{align*}
  \bm a^\top(A\circ B)\,\bm a
  &=\sum_{i,j}a_ia_jA_{ij}\sum_m\lambda_mv_{mi}v_{mj}
   =\sum_m\lambda_m\sum_{i,j}(a_iv_{mi})\,A_{ij}\,(a_jv_{mj})\\
  &=\sum_m\lambda_m\,\bm u_m^\top A\,\bm u_m\ \ge\ 0,
\end{align*}
em que $\bm u_m=(a_1v_{m1},\dots,a_nv_{mn})$, e cada parcela é não negativa porque
$\lambda_m\ge0$ e $A$ é positiva semidefinida. Como isso vale para todo $n$ e quaisquer
pontos, o teorema dá que $k_1k_2$ é um \emph{kernel}.

\smallskip
\textbf{(b)} A função $h(\x)h(\x')k(\x,\x')$ é simétrica, e a sua matriz nos pontos
$\x_1,\dots,\x_n$ é $DKD$, com $K=\big[k(\x_i,\x_j)\big]$ e
$D=\operatorname{diag}\big(h(\x_1),\dots,h(\x_n)\big)$. Para $\bm a\in\R^n$,
$\bm a^\top DKD\,\bm a=(D\bm a)^\top K(D\bm a)\ge0$. (Pela definição, é o mesmo que tomar a
função $h\varphi$ no lugar de $\varphi$.)

Para a exponencial, $e^{k(\x,\x')}=\sum_{m\ge0}k(\x,\x')^m/m!$. A potência $k^0$ é a
constante $1$, que é um \emph{kernel}, e, se $k^m$ é um \emph{kernel}, $k^{m+1}=k^m\cdot k$
também é, pelo item (a); por indução, todas as potências são \emph{kernels}. As somas
parciais $S_M=\sum_{m=0}^Mk^m/m!$ são, então, \emph{kernels}, como somas de \emph{kernels}
multiplicados por constantes positivas. Fixados os pontos e $\bm a\in\R^n$, cada número
$\bm a^\top\big[S_M(\x_i,\x_j)\big]\bm a$ é não negativo, e, quando $M\to\infty$, converge para
$\bm a^\top\big[e^{k(\x_i,\x_j)}\big]\bm a$, porque é uma soma finita de termos que convergem.
Um limite de números não negativos é não negativo. A matriz de $e^k$ é, portanto, positiva
semidefinida, e $e^k$, que é simétrica, é um \emph{kernel}.

\smallskip
\textbf{(c)} Como $\norm{\x-\x'}^2=\norm{\x}^2-2\inner{\x}{\x'}+\norm{\x'}^2$,
\[
  e^{-\gamma\norm{\x-\x'}^2}
  =e^{-\gamma\norm{\x}^2}\,e^{-\gamma\norm{\x'}^2}\,e^{2\gamma\inner{\x}{\x'}} .
\]
O produto interno é um \emph{kernel}, e $2\gamma\inner{\x}{\x'}$ também, porque $2\gamma>0$.
Pelo item (b), $e^{2\gamma\inner{\x}{\x'}}$ é um \emph{kernel}, e, ainda pelo item (b), com
$h(\x)=e^{-\gamma\norm{\x}^2}$, o produto acima é um \emph{kernel}. O argumento vale para todo
$p$ e todo $\gamma>0$, e não precisa exibir $\varphi$; com $p=1$, expandir
$e^{2\gamma xx'}$ em série reproduz a função $\varphi$ do exemplo das notas.

\smallskip
\textbf{(d)} Pela bilinearidade e pela simetria do produto interno de $\mathcal H$,
\begin{align*}
  \norm{\varphi(\x)-\varphi(\x')}^2
  &=\inner{\varphi(\x)}{\varphi(\x)}-2\inner{\varphi(\x)}{\varphi(\x')}+\inner{\varphi(\x')}{\varphi(\x')}\\
  &=k(\x,\x)-2\,k(\x,\x')+k(\x',\x') .
\end{align*}
No \emph{kernel} radial, $k(\x,\x)=e^0=1$, de modo que $\norm{\varphi(\x)}=1$ para todo $\x$,
e
\[
  \norm{\varphi(\x)-\varphi(\x')}^2=2\big(1-e^{-\gamma\norm{\x-\x'}^2}\big) .
\]
Essa quantidade é nula se e somente se $\x=\x'$, cresce com $\norm{\x-\x'}$ e é sempre menor
que $2$: os pontos vão para a esfera unitária de $\mathcal H$, a distâncias menores que
$\sqrt2$, isto é, com produto interno $k(\x,\x')>0$ e ângulo menor que $90^\circ$.

Dois pontos ficam à distância $1$ em $\mathcal H$ quando
$2\big(1-e^{-\gamma\norm{\x-\x'}^2}\big)=1$, isto é, quando $k(\x,\x')=1/2$, a meio caminho
entre pontos iguais ($k=1$) e ortogonais ($k=0$). Isso acontece se e somente se
\[
  \norm{\x-\x'}=\sqrt{\frac{\log2}{\gamma}} ,
\]
que vale $\sqrt{2\log2}\approx1{,}177$ com $\gamma=0{,}5$ e $\sqrt{2\log2}/10\approx0{,}118$
com $\gamma=50$, dez vezes menos.

A distância em $\mathcal H$ depende de $\x$ e $\x'$ só por meio de
$\sqrt\gamma\,\norm{\x-\x'}$: multiplicar $\gamma$ por $100$ tem o mesmo efeito que
multiplicar todas as covariáveis por $10$, e $1/\sqrt\gamma$ faz o papel de uma unidade de
comprimento em $\R^p$. É também a escala em que age cada parcela da função de decisão
$f(\x)=\beta_0+\sum_{i\in\mathcal S}\alpha_ik(\x,\x_i)$: $k(\x,\x_i)$ cai à metade quando
$\norm{\x-\x_i}=\sqrt{\log2/\gamma}$ e decai depressa a partir daí. Com $\gamma=50$, a
influência de cada vetor de suporte se restringe a uma vizinhança de raio da ordem de
$0{,}1$ em torno dele, a fronteira pode contornar observações isoladas, e o ajuste fica
flexível, com risco de superajuste. Com $\gamma=0{,}5$, o raio é dez vezes maior, muitos
vetores de suporte pesam em cada ponto, e a fronteira é suave.
\end{solucao}


\begin{criterio}
\textbf{(a)} 0,2 pela simetria de $A\circ B$; 0,4 pela decomposição espectral de $B$,
com autovalores não negativos porque $B$ é positiva semidefinida; 0,6 por reescrever
$\bm a^\top(A\circ B)\,\bm a$ como $\sum_m\lambda_m\,\bm u_m^\top A\,\bm u_m$; 0,4 pela
conclusão, parcela a parcela, e pelo teorema, porque a conta vale para todo $n$ e
quaisquer pontos. Aceitar integralmente uma prova correta por outro caminho, como a
construção de uma função $\varphi$ para o produto. Citar o teorema do produto de Schur
sem demonstrá-lo vale 0,4 do item.

\textbf{(b)} 0,5 por $h(\x)h(\x')k(\x,\x')$: a simetria e a matriz $DKD$, com
$\bm a^\top DKD\,\bm a=(D\bm a)^\top K(D\bm a)\ge0$, ou a função $h\varphi$ no lugar de
$\varphi$. 0,8 pela exponencial: 0,3 pelas potências, por indução e pelo item (a); 0,2
pelas somas parciais, combinações de \emph{kernels} com coeficientes positivos; 0,3 pela
passagem ao limite, com os pontos e $\bm a$ fixados. Afirmar que um limite de
\emph{kernels} é um \emph{kernel} sem dizer em que sentido vale 0,1 desses 0,3.

\textbf{(c)} 0,3 pela fatoração
$e^{-\gamma\norm{\x-\x'}^2}=e^{-\gamma\norm{\x}^2}e^{-\gamma\norm{\x'}^2}e^{2\gamma\inner{\x}{\x'}}$;
0,4 por $e^{2\gamma\inner{\x}{\x'}}$ ser um \emph{kernel}, porque $2\gamma\inner{\x}{\x'}$ é
um múltiplo positivo do produto interno e pelo item (b); 0,4 pelo item (b) com
$h(\x)=e^{-\gamma\norm{\x}^2}$ e pela observação de que o argumento não depende de $p$
nem do valor de $\gamma>0$. Usar os itens (a) e (b) sem tê-los demonstrado não desconta
neste item.
\end{criterio}

%==============================================================================
\questao{}{4,0 pontos}
%==============================================================================
% Aula 10 --- uma divisão em que o Gini e a entropia discordam, e a forma da curva que explica isso.
Ana é bióloga e usa uma árvore de classificação para distinguir mudas de duas espécies, a
espécie $1$ e a espécie $0$, a partir de características medidas no viveiro. Num nó com $12$
mudas, $6$ de cada espécie, a árvore compara duas divisões:
\begin{description}[leftmargin=1.4cm,style=nextline]
  \item[A] pela altura, separa dois grupos de $6$ mudas, um com $2$ da espécie $1$ e $4$ da
        espécie $0$ e outro com $4$ da espécie $1$ e $2$ da espécie $0$;
  \item[B] por uma mancha rara na folha, isola uma única muda, da espécie $0$, das outras $11$,
        que têm $6$ da espécie $1$ e $5$ da espécie $0$.
\end{description}
Numa região com proporção $p$ de mudas da espécie $1$, o índice de Gini é $G(p)=2p(1-p)$, a
entropia é $D(p)=-p\log_2p-(1-p)\log_2(1-p)$, com $0\log_20=0$, e o erro de classificação é
$E(p)=\min(p,1-p)$. Uma divisão em regiões com $n_1$ e $n_2$ mudas e proporções $p_1$ e $p_2$
tem impureza ponderada $\big[n_1I(p_1)+n_2I(p_2)\big]/(n_1+n_2)$, para cada medida $I$, e a
árvore escolhe a divisão de menor impureza ponderada.
\begin{itens}
  \item \pts{2,0} Calcule a impureza ponderada das divisões A e B pelas três medidas e diga
        qual divisão cada medida escolhe. \emph{Dados:} $\log_23\approx1{,}5850$,
        $\log_25\approx2{,}3219$ e $\log_211\approx3{,}4594$.
  \item \pts{2,0} Seja $I$ uma medida de impureza qualquer para duas classes, com
        $I(p)=I(1-p)$, $I(0)=0$ e $I(6/11)>0$. Mostre que A tem impureza ponderada
        menor que B se e somente se $I(1/3)/I(6/11)<11/12$. Calcule essa razão para as
        três medidas e explique que propriedade da curva da entropia a leva a escolher B.
\end{itens}

\begin{solucao}
\textbf{(a)} Na divisão A, as proporções da espécie $1$ são $1/3$ e $2/3$, cada uma num filho
com metade das mudas. Na divisão B, são $0$, numa muda, e $6/11$, nas outras $11$. Pelo Gini,
\[
  \text{A: }\ \tfrac12\cdot2\cdot\tfrac13\cdot\tfrac23+\tfrac12\cdot2\cdot\tfrac23\cdot\tfrac13=\tfrac49\approx0{,}444,
  \qquad
  \text{B: }\ \tfrac1{12}\cdot0+\tfrac{11}{12}\cdot2\cdot\tfrac6{11}\cdot\tfrac5{11}=\tfrac5{11}\approx0{,}455 .
\]
Pela entropia, como $D(1/3)=D(2/3)=\tfrac13\log_23+\tfrac23\log_2\tfrac32=\log_23-\tfrac23$,
\begin{align*}
  \text{A: }\ &\log_23-\tfrac23\approx0{,}918,\\
  \text{B: }\ &\tfrac{11}{12}D\big(\tfrac6{11}\big)
  =\tfrac1{12}\Big(6\log_2\tfrac{11}6+5\log_2\tfrac{11}5\Big)
  =\tfrac1{12}\big(11\log_211-6\log_26-5\log_25\big)\approx0{,}911 ,
\end{align*}
com $\log_26=1+\log_23$. Pelo erro de classificação, A erra $2$ mudas em cada filho, $4/12=1/3$,
e B erra as $5$ mudas da espécie $0$ no filho grande, $5/12\approx0{,}417$. O Gini e o erro
escolhem A; a entropia escolhe B. O nó tem $G=0{,}5$ e $D=1$, e as reduções correspondentes são
$0{,}056$ (A) e $0{,}045$ (B) no Gini, e $0{,}082$ (A) e $0{,}089$ (B) na entropia.

\smallskip
\textbf{(b)} Pela simetria, a impureza ponderada de A é
$\tfrac12I(1/3)+\tfrac12I(2/3)=I(1/3)$, e a de B é
$\tfrac1{12}I(0)+\tfrac{11}{12}I(6/11)=\tfrac{11}{12}I(6/11)$. Como $I(6/11)>0$,
\[
  I\big(\tfrac13\big)<\tfrac{11}{12}\,I\big(\tfrac6{11}\big)
  \iff \frac{I(1/3)}{I(6/11)}<\frac{11}{12}\approx0{,}917 .
\]
As razões são $(4/9)/(60/121)=121/135\approx0{,}896$ para o Gini,
$0{,}918/0{,}994\approx0{,}924$ para a entropia e $(1/3)/(5/11)=11/15\approx0{,}733$ para o erro: só
a entropia passa de $11/12$.

A comparação tem uma leitura simples. A divisão A leva os dois filhos da proporção $1/2$ à
proporção $1/3$; a divisão B torna pura $1/12$ das mudas e deixa as outras numa região quase tão
misturada quanto o nó, porque $6/11\approx0{,}545$ está perto de $1/2$. Com
$I(6/11)\approx I(1/2)$, a pergunta é, em essência, se $I(1/3)$ fica abaixo de $11/12$ do
máximo: $G(1/3)/G(1/2)=8/9\approx0{,}889$, mas $D(1/3)/D(1/2)\approx0{,}918$. A entropia é mais
achatada perto do máximo: trata uma região com um terço de uma espécie como quase tão impura
quanto uma equilibrada, e por isso dá pouco valor à purificação parcial de A e, em comparação,
mais valor à purificação completa de um pedaço pequeno, que é o que B faz.
\end{solucao}


\begin{criterio}
\textbf{(a)} 0,5 pelo Gini, $4/9\approx0{,}444$ em A e $5/11\approx0{,}455$ em B; 0,8 pela
entropia, 0,4 por divisão, $\log_23-2/3\approx0{,}918$ em A e cerca de $0{,}911$ em B; 0,4
pelo erro de classificação, $1/3$ em A e $5/12$ em B; 0,3 pelas escolhas: o Gini e o erro
escolhem A, e a entropia escolhe B. Aceitar desvios de até $0{,}002$ nos valores da
entropia, que dependem dos logaritmos dados, e impurezas sem a divisão por $12$, que
não mudam a comparação. Uma conta errada que inverta a escolha da entropia perde os
pontos dessa conta e dessa escolha, mas não os das outras medidas.

\textbf{(b)} 0,5 por escrever as impurezas ponderadas como $I(1/3)$ e
$\tfrac{11}{12}I(6/11)$, usando a simetria e $I(0)=0$; 0,3 pela equivalência, dividindo
por $I(6/11)>0$; 0,6 pelas três razões, 0,2 cada ($121/135\approx0{,}896$ no Gini, cerca de
$0{,}924$ na entropia e $11/15\approx0{,}733$ no erro), com a conclusão de que só a da
entropia passa de $11/12$; 0,6 pela explicação: perto de $p=1/2$ a entropia é mais
achatada que o Gini ($D(1/3)/D(1/2)\approx0{,}918$, contra $G(1/3)/G(1/2)=8/9$), e por isso
dá pouco valor à purificação parcial de A diante da purificação completa de uma única
muda em B. Dizer apenas que as curvas têm formas diferentes, sem ligar a forma aos
valores em $1/3$ e perto de $1/2$, vale 0,2 desses 0,6.
\end{criterio}

\vfill
\begin{center}\small\itshape Boa prova!\end{center}

\end{document}
